\documentclass[11pt,a4paper]{article}
\usepackage[utf8]{inputenc}
\usepackage[T1]{fontenc}
\usepackage{lmodern}
\usepackage[margin=2.2cm]{geometry}
\usepackage{amsmath,amssymb,amsthm,mathtools}
\usepackage{enumitem}
\usepackage{xcolor}
\usepackage[most]{tcolorbox}
\usepackage{tikz}
\usetikzlibrary{arrows.meta,calc}
\usepackage{fancyhdr}
\usepackage[hidelinks]{hyperref}
\usepackage{multicol}

\definecolor{azul}{RGB}{40,70,150}
\definecolor{verde}{RGB}{30,120,70}
\definecolor{rojo}{RGB}{170,40,40}
\definecolor{gris}{RGB}{90,90,90}

\setlength{\parindent}{0pt}
\setlength{\parskip}{5pt}

\newtcolorbox{resumen}[1][]{colback=azul!5,colframe=azul,fonttitle=\bfseries,title={#1},breakable,arc=2pt,boxrule=0.8pt}
\newtcolorbox{receta}[1][Receta]{colback=verde!6,colframe=verde,fonttitle=\bfseries,title={#1},breakable,arc=2pt,boxrule=0.8pt}
\newtcolorbox{ojo}[1][Errores típicos]{colback=rojo!5,colframe=rojo,fonttitle=\bfseries,title={#1},breakable,arc=2pt,boxrule=0.8pt}
\newtcolorbox{ejemplo}[1][]{colback=white,colframe=gris,fonttitle=\bfseries,title={#1},breakable,arc=2pt,boxrule=0.6pt}

\newcounter{ejer}
\newcommand{\ejercicio}[2][]{\refstepcounter{ejer}\par\medskip\noindent\textbf{Ejercicio \theejer.}\ \textit{\small #2}\ifx&#1&\else\hfill{\small\textcolor{gris}{[#1]}}\fi\par}
\newcommand{\sol}[1]{\par\medskip\noindent\textbf{Ejercicio #1.}\ }

\newcommand{\R}{\mathbb{R}}
\newcommand{\dA}{\,dA}
\newcommand{\dV}{\,dV}

\pagestyle{fancy}
\fancyhf{}
\renewcommand{\headrulewidth}{0.4pt}
\fancyfoot[C]{\small\thepage}
\renewcommand{\figurename}{Figura}
\renewcommand{\tablename}{Tabla}
\fancyhead[L]{\small Cálculo para Ciencia de Datos -- IMT2220}
\fancyhead[R]{\small Guía 4: Integrales triples}

\begin{document}

\begin{center}
{\LARGE\bfseries Guía 4: Integrales triples,\\[2pt] coordenadas cilíndricas y esféricas}\\[6pt]
{\large Cálculo para Ciencia de Datos (IMT2220) -- Preparación Examen 2026-2}\\[3pt]
{\small Material: Clases 21 y 22 \quad$\cdot$\quad Ayudantías 9 y 10 \quad$\cdot$\quad Exámenes 2022-1 y 2023-1}
\end{center}

\section{Resumen}

\begin{resumen}[Integral triple y Fubini]
Se define con sumas de Riemann sobre cajas pequeñas $\Delta V=\Delta x\,\Delta y\,\Delta z$. Si $f$ es continua en la caja $B$, vale \textbf{Fubini en cualquiera de los 6 órdenes}. Con $f=1$ se obtiene el volumen; con $f$ densidad de masa, la masa; con $f$ densidad de probabilidad, la probabilidad.
\end{resumen}

\begin{resumen}[Regiones elementales]
\textbf{Respecto al plano $xy$:} $T=\{(x,y)\in R,\ k_1(x,y)\le z\le k_2(x,y)\}$, y
\[
\iiint_T f\dV=\iint_R\Big(\int_{k_1(x,y)}^{k_2(x,y)}f\,dz\Big)dA .
\]
Igual respecto a $yz$ (se integra primero $x$) o a $xz$ (primero $y$).

\textbf{Receta:} proyectar el sólido sobre un plano para obtener $R$; para cada punto de $R$, ver entre qué dos superficies se mueve la variable restante. Los límites de cada integral dependen solo de las variables que \emph{todavía no} se integran.
Conviene elegir el plano según la geometría: si el eje de un cilindro es el eje $y$, se proyecta sobre $xz$.
\end{resumen}

\begin{resumen}[Cambio de variables en $\R^3$]
\[
\iiint_D f\,dx\,dy\,dz=\iiint_{D^*}f\big(T(u,v,w)\big)\left|\frac{\partial(x,y,z)}{\partial(u,v,w)}\right|du\,dv\,dw .
\]
Igual que en 2D, se puede usar el jacobiano de la inversa: $\frac{\partial(x,y,z)}{\partial(u,v,w)}=1\big/\frac{\partial(u,v,w)}{\partial(x,y,z)}$.
\end{resumen}

\begin{resumen}[Cilíndricas y esféricas]
\begin{minipage}[t]{0.47\textwidth}
\textbf{Cilíndricas} (simetría respecto a un eje):
\[
x=r\cos\theta,\ y=r\sin\theta,\ z=z
\]
\[
\boxed{dV=r\,dr\,d\theta\,dz}
\]
Para cilindros, paraboloides ($z=r^2$) y conos ($z=r$). Son polares en el plano más la altura.
\end{minipage}\hfill
\begin{minipage}[t]{0.49\textwidth}
\textbf{Esféricas} (simetría respecto a un punto):
\[
x=\rho\sin\varphi\cos\theta,\ y=\rho\sin\varphi\sin\theta,\ z=\rho\cos\varphi
\]
\[
\boxed{dV=\rho^2\sin\varphi\,d\rho\,d\varphi\,d\theta}
\]
$\rho\ge0$, \ $0\le\varphi\le\pi$ (\textbf{desde el eje $z$ positivo}), \ $0\le\theta\le2\pi$. Además $\rho^2=x^2+y^2+z^2$ y $r=\rho\sin\varphi$.
\end{minipage}

\medskip
\textbf{Superficies en esféricas:} $\rho=c$ es una esfera, $\varphi=c$ un cono y $\theta=c$ un semiplano.\\
$x^2+y^2+(z-1)^2=1\iff\rho=2\cos\varphi$; \quad $z=\sqrt{x^2+y^2}\iff\varphi=\frac\pi4$; \quad $z=\sqrt{3(x^2+y^2)}\iff\varphi=\frac\pi6$; \quad $z=1\iff\rho=\frac1{\cos\varphi}$.

\textbf{Esféricas adaptadas al elipsoide} $\frac{x^2}{a^2}+\frac{y^2}{b^2}+\frac{z^2}{c^2}\le1$: $x=a\rho\sin\varphi\cos\theta$, etc., con jacobiano $abc\,\rho^2\sin\varphi$. Volumen $\frac43\pi abc$.
\end{resumen}

\begin{resumen}[Probabilidad en $\R^3$]
$p\ge0$, $\iiint_{\R^3}p=1$ y $P((X,Y,Z)\in E)=\iiint_E p\dV$. \textbf{Independencia:} $p(x,y,z)=p_X(x)p_Y(y)p_Z(z)$.
Uniforme en la bola $B_R$: $P(\rho\le r)=(r/R)^3$. En $\R^n$ es $(r/R)^n$, y el volumen se concentra en la cáscara (maldición de la dimensionalidad).
\end{resumen}

\begin{ojo}
\begin{itemize}[leftmargin=*,itemsep=1pt]
\item Olvidar $r$ (cilíndricas) o $\rho^2\sin\varphi$ (esféricas).
\item Confundir $\varphi$ (desde el eje $z$, de $0$ a $\pi$) con $\theta$ (en el plano $xy$, de $0$ a $2\pi$).
\item Poner en los límites variables que ya se integraron.
\item No dibujar la intersección de las superficies: ahí están los límites (por ejemplo, el cono $z=r$ y el plano $z=2$ se cortan en $r=2$).
\end{itemize}
\end{ojo}

\section{Ejemplos resueltos}

\begin{ejemplo}[Ejemplo 1 (Clase 21). Volumen bajo un paraboloide]
Sólido bajo $z=4-x^2-y^2$ y sobre $z=0$. Proyección: el disco $r\le2$. En cilíndricas, $0\le z\le4-r^2$:
\[
V=\int_0^{2\pi}\!\!\int_0^2\!\!\int_0^{4-r^2}r\,dz\,dr\,d\theta=2\pi\int_0^2(4-r^2)r\,dr=2\pi\Big[2r^2-\frac{r^4}4\Big]_0^2=8\pi .
\]
\end{ejemplo}

\begin{ejemplo}[Ejemplo 2 (Clase 21). Cono cerrado por un plano]
$\iiint_\Omega(1+(x^2+y^2)^2)\dV$, con $\Omega$ entre el cono $z=\sqrt{x^2+y^2}$ y el plano $z=2$. Se cortan en $r=2$, así que $\Omega: r\le z\le2$, $0\le r\le2$:
\[
\int_0^{2\pi}\!\!\int_0^2(2-r)(1+r^4)\,r\,dr\,d\theta=2\pi\Big(4-\frac83+\frac{64}3-\frac{128}7\Big)=\frac{184\pi}{21}.
\]
\end{ejemplo}

\begin{ejemplo}[Ejemplo 3 (Clase 22). Cono con tapa esférica]
$\iiint_D z\dV$, con $D$ entre el cono $z=\sqrt{x^2+y^2}$ y la esfera $x^2+y^2+z^2=z$. En esféricas, la esfera es $\rho=\cos\varphi$ y el cono es $\varphi=\frac\pi4$:
\[
\int_0^{2\pi}\!\!\int_0^{\pi/4}\!\!\int_0^{\cos\varphi}(\rho\cos\varphi)\rho^2\sin\varphi\,d\rho\,d\varphi\,d\theta
=\frac{2\pi}4\int_0^{\pi/4}\cos^5\varphi\sin\varphi\,d\varphi=\frac\pi{12}\Big(1-\frac18\Big)=\frac{7\pi}{96}.
\]
\end{ejemplo}

\begin{ejemplo}[Ejemplo 4 (Clase 21). Paralelepípedo inclinado]
$E=\{0\le x+y\le1,\ 0\le y+z\le1,\ 0\le x+z\le1\}$. Con $u=x+y$, $v=y+z$, $w=x+z$ se tiene $(u,v,w)\in[0,1]^3$ y
$\frac{\partial(u,v,w)}{\partial(x,y,z)}=\det\begin{psmallmatrix}1&1&0\\0&1&1\\1&0&1\end{psmallmatrix}=2$. Entonces $\mathrm{Vol}(E)=\iiint_{[0,1]^3}\frac12\,du\,dv\,dw=\frac12$.
\end{ejemplo}

\section{Ejercicios}

\ejercicio[Clase 21]{Fubini en una caja}
Calcule $\iiint_B(x^2y+zy^2)\dV$, con $B=[0,1]\times[1,2]\times[-1,1]$.

\ejercicio[Ay 10]{Región elemental en dos órdenes}
Sea $E$ el sólido limitado por el cilindro parabólico $y=x^2$ y los planos $z=0$ e $y+z=1$.
\begin{enumerate}[label=(\alph*),itemsep=0pt]
\item Describa $E$ respecto al plano $xy$ y calcule su volumen.
\item Calcule $\iiint_E z\dV$.
\end{enumerate}

\ejercicio[Clase 21]{Cilíndricas}
\begin{enumerate}[label=(\alph*),itemsep=0pt]
\item Volumen del sólido dentro del cilindro $x^2+y^2=4$, entre $z=0$ y $z=4-y$.
\item $\iiint_E\sqrt{x^2+y^2}\dV$ con $E=\{x^2+y^2\le1,\ 0\le z\le x^2+y^2\}$.
\end{enumerate}

\ejercicio[Examen 2022-1]{Cilindro cortado por un plano}
Sea $A$ la región dentro del cilindro $x^2+y^2=9$, sobre el plano $z=0$ y bajo el plano $y+z=5$. Dibuje $A$ y calcule $\iiint_A\sqrt{x^2+y^2}\dV$.

\ejercicio[Examen 2023-1]{Cilíndricas y esféricas}
\begin{enumerate}[label=(\alph*),itemsep=0pt]
\item Calcule $\iiint_A\sqrt{x^2+y^2}\dV$, con $A$ la región dentro del cilindro $x^2+y^2=1$, bajo el plano $z=4$ y sobre el paraboloide $z=1-x^2-y^2$.
\item Calcule $\iiint_B e^{(x^2+y^2+z^2)^{3/2}}\dV$, con $B$ la bola unitaria centrada en el origen.
\end{enumerate}

\ejercicio[Ay 9]{Llenar un recipiente}
Un recipiente tiene la forma del sólido limitado por el paraboloide $z=x^2+y^2$ y el plano $z=4$.
\begin{enumerate}[label=(\alph*),itemsep=0pt]
\item Describa el sólido en cilíndricas y calcule su volumen.
\item Si se llena hasta la altura $h$ ($0<h\le4$), calcule el volumen de líquido $V(h)$. ¿Hasta qué altura hay que llenarlo para que contenga la mitad? ¿Por qué no es $h=2$?
\end{enumerate}

\ejercicio[Ay 10]{Esfera perforada}
Se perfora una esfera de radio 2 con un taladro de radio 1 a lo largo del eje $z$: $E=\{x^2+y^2+z^2\le4,\ x^2+y^2\ge1\}$. Describa $E$ en cilíndricas y calcule su volumen. ¿Cuánto material se removió?

\ejercicio[Clase 22]{Esféricas}
\begin{enumerate}[label=(\alph*),itemsep=0pt]
\item Escriba en esféricas la esfera $x^2+y^2+z^2=4z$ y el cono $z=\sqrt{3(x^2+y^2)}$.
\item Calcule el volumen del sólido sobre el cono $\varphi=\frac\pi4$ y dentro de la esfera $\rho=2$.
\item Calcule $\iiint_{B_1}(x^2+y^2+z^2)\dV$ sobre la bola unitaria.
\item Calcule el volumen del elipsoide $\frac{x^2}4+\frac{y^2}9+z^2\le1$ usando esféricas adaptadas.
\end{enumerate}

\ejercicio[Ay 10]{Balón de rugby}
El balón es el elipsoide $\frac{x^2}{16}+\frac{y^2}9+\frac{z^2}4\le1$, con densidad $\rho_m=1-\big(\frac{x^2}{16}+\frac{y^2}9+\frac{z^2}4\big)$. Calcule su volumen y su masa.

\ejercicio[Ay 10]{Pisapapeles}
Un tetraedro está limitado por los planos coordenados y el plano $x+y+z=2$ (en cm), con densidad $\rho_m=1+z$ (g/cm$^3$). Calcule su masa integrando al final en $z$ (corte a altura $z$ de área $A(z)$) y compárela con la de un vidrio homogéneo de densidad 1.

\ejercicio[Clase 22]{Tres micros}
Tres micros pasan cada 10 minutos, de forma independiente. Los tiempos de espera $T_1,T_2,T_3$ son uniformes en $[0,10]$. Calcule $P(T_1+T_2+T_3\le5)$.

\newpage
\section{Soluciones}

\sol{1}
$\int_{-1}^1\int_1^2\int_0^1(x^2y+zy^2)\,dx\,dy\,dz=\int_{-1}^1\Big(\frac12+\frac73z\Big)dz=1$.

\sol{2}
(a) La proyección es $R=\{-1\le x\le1,\ x^2\le y\le1\}$ y $0\le z\le1-y$.
$V=\int_{-1}^1\int_{x^2}^1(1-y)\,dy\,dx=\int_{-1}^1\frac{(1-x^2)^2}2dx=\frac8{15}$.
(b) $\int_{-1}^1\int_{x^2}^1\frac{(1-y)^2}2\,dy\,dx=\int_{-1}^1\frac{(1-x^2)^3}6dx=\frac{16}{105}$.

\sol{3}
(a) $\int_0^{2\pi}\int_0^2(4-r\sin\theta)\,r\,dr\,d\theta=16\pi$ (el término con $\sin\theta$ se anula).
(b) $\int_0^{2\pi}\int_0^1\int_0^{r^2}r\cdot r\,dz\,dr\,d\theta=2\pi\int_0^1r^4dr=\frac{2\pi}5$.

\sol{4}
En cilíndricas: $0\le r\le3$, $0\le\theta\le2\pi$, $0\le z\le5-r\sin\theta$.
$\int_0^{2\pi}\int_0^3 r(5-r\sin\theta)\,r\,dr\,d\theta=5\cdot9\cdot2\pi-0=90\pi$.

\sol{5}
(a) $0\le r\le1$, $1-r^2\le z\le4$: $\int_0^{2\pi}\int_0^1 r(3+r^2)\,r\,dr\,d\theta=2\pi\big(1+\frac15\big)=\frac{12\pi}5$.
(b) En esféricas: $4\pi\int_0^1e^{\rho^3}\rho^2\,d\rho=4\pi\cdot\frac{e-1}3=\frac{4\pi(e-1)}3$.

\sol{6}
(a) $r^2\le z\le4$, $0\le r\le2$: $V=\int_0^{2\pi}\int_0^2(4-r^2)r\,dr\,d\theta=8\pi$.
(b) $V(h)=\int_0^{2\pi}\int_0^{\sqrt h}(h-r^2)r\,dr\,d\theta=\frac{\pi h^2}2$. La mitad es $\frac{\pi h^2}2=4\pi$, es decir $h=2\sqrt2\approx2.83$. No es $h=2$ porque el recipiente es más ancho arriba.

\sol{7}
$1\le r\le2$ y $-\sqrt{4-r^2}\le z\le\sqrt{4-r^2}$: $V=2\pi\int_1^2 2r\sqrt{4-r^2}\,dr=2\pi\cdot\frac23\cdot3^{3/2}=4\sqrt3\,\pi\approx21.8$.
Removido: $\frac{32\pi}3-4\sqrt3\pi\approx11.7$.

\sol{8}
(a) $\rho=4\cos\varphi$; \ $\varphi=\frac\pi6$.
(b) $\int_0^{2\pi}\int_0^{\pi/4}\int_0^2\rho^2\sin\varphi\,d\rho\,d\varphi\,d\theta=\frac{16\pi}3\Big(1-\frac{\sqrt2}2\Big)\approx4.91$.
(c) $\int_0^{2\pi}\int_0^\pi\int_0^1\rho^4\sin\varphi\,d\rho\,d\varphi\,d\theta=\frac{4\pi}5$.
(d) $a=2$, $b=3$, $c=1$: $\frac43\pi\cdot6=8\pi$.

\sol{9}
Con esféricas adaptadas ($a,b,c=4,3,2$, jacobiano $24\rho^2\sin\varphi$) el balón es $\rho\le1$ y la densidad es $1-\rho^2$.
Volumen $=\frac43\pi\cdot24=32\pi$. Masa $=24\cdot4\pi\int_0^1(1-\rho^2)\rho^2d\rho=96\pi\Big(\frac13-\frac15\Big)=\frac{64\pi}5$.

\sol{10}
El corte a altura $z$ es el triángulo $x,y\ge0$, $x+y\le2-z$, de área $A(z)=\frac{(2-z)^2}2$. Entonces
$M=\int_0^2(1+z)\frac{(2-z)^2}2\,dz=2$ g. Homogéneo: masa $=$ volumen $=\frac{2^3}6=\frac43$ g.

\sol{11}
La región es el tetraedro $t_i\ge0$, $t_1+t_2+t_3\le5$, de volumen $\frac{5^3}6$, y la densidad es $\frac1{1000}$: $P=\frac{125/6}{1000}=\frac1{48}\approx0.021$.

\end{document}
